\documentclass[12pt]{article}
\usepackage[a4paper,margin=1in]{geometry}
\usepackage[utf8]{inputenc}
\usepackage[T1]{fontenc}
\usepackage{amsmath,amssymb}
\usepackage{textcomp}
\usepackage{parskip}
\begin{document}
\begin{center}
{\Large \textbf{2022 STEP II}}
\end{center}

\subsection*{Page 1}
\begin{flushleft}
\bigskip\textbf{Section A: Pure Mathematics}\par
\bigskip\textbf{Question 1.} (i) By integrating one of the two terms in the integrand by parts, or otherwise, find\par
\[\leq \geq )\]
( 3\par
3x\par
\[2 1 + x^3 + \sqrt\{dx\} .\]
\[1 + x^3\]
\textbf{(ii)} Find \leq\par
sin x\par
(x + 2) dx .\par
\[x^3\]
x\par
e\par
\textbf{(iii)} (a) Sketch the graph with equation y = , giving the coordinates of any stationary\par
x\par
points.\par
(b) Find a if\par
\[\leq \leq\]
2a x 2a x\par
e e\par
\[dx = dx .\]
\[x x^2\]
a a\par
(c) Show that it is not possible to find distinct integers m and n such that\par
\[\leq \leq\]
n x n x\par
e e\par
\[dx = dx .\]
\[x x^2\]
m m\par
\bigskip\textbf{Question 2.} A sequence un, where n = 1, 2, . . ., is said to have degree d if un, as a function of n, is a\par
polynomial of degree d.\par
\textbf{(i)} Show that, in any sequence un (n = 1, 2, . . .) that satisfies un+1 =\par
1 (un+2 + un) for\par
\[all n \leq 1, there is a constant difference between successive terms.\]
\[Deduce that any sequence un for which un+1 = 1\]
\[(un+2 + un), for all n \leq 1, has degree\]
at most 1.\par
\textbf{(ii)} The sequence vn (n = 1, 2, . . .) satisfies vn+1 = 1\par
\[(vn+2 + vn) - p for all n \leq 1, where\]
\[p is a non-zero constant. By writing vn = tn + pn , show that the sequence vn has\]
degree 2.\par
\[Given that v^1 = v^2 = 0, find vn in terms of n and p.\]
\textbf{(iii)} The sequence wn (n = 1, 2, . . .) satisfies wn+1 = 1\par
\[(wn+2 + wn) - an - b for all n \leq 1,\]
where a and b are constants with a \ne 0. Show that the sequence wn has degree 3.\par
\[Given that w^1 = w^2 = 0, find wn in terms of n, a and b.\]
\end{flushleft}
\subsection*{Page 2}
\begin{flushleft}
\bigskip\textbf{Question 3.} The Fibonacci numbers are defined by F0 = 0, F1 = 1 and, for n \leq 0, Fn+2 = Fn+1 + Fn.\par
\textbf{(i)} Prove that Fr \geq 2 r-n Fn for all n \leq 1 and all r \leq n.\par
\[\intn\]
\textbf{(ii)} Let Sn =\par
F\par
r\par
r\par
.\par
\[r=1\]
Show that\par
\[\intn \intn \intn\]
F\par
r\par
r\par
+\par
-\par
-\par
F\par
r\par
r\par
-1\par
-\par
F\par
r\par
r\par
-\par
-\par
\[= 89Sn - 10F1 - F^0 +\]
F\par
n\par
n\par
+\par
F\par
n\par
n\par
+\par
-\par
.\par
\[r=1 r=1 r=1\]
\[\int\infty \int\infty\]
\textbf{(iii)} Show that\par
Fr\par
\[= 10 and that\]
Fr\par
\[< 2 \times 10 -6 . Hence find, with justification,\]
10r 89 10r\par
\[r=1 r=7\]
the first six digits after the decimal point in the decimal expansion of .\par
r\par
\textbf{(iv)} Find, with justification, a number of the form with r and s both positive integers\par
s\par
less than 10000 whose decimal expansion starts\par
0.0001010203050813213455 . . . .\par
\end{flushleft}
\subsection*{Page 3}
\begin{flushleft}
\[\sum \sum \sum \sum\]
\bigskip\textbf{Question 4.} (i) Show that the function f, given by the single formula f(x) = \sum x \sum - \sum x - 5 \sum + 1, can be\par
written without using modulus signs as\par
⎧\par
\[⎨ -4 x \geq 0 ,\]
\[f(x) = 2x - 4 0 \geq x \geq 5 ,\]
⎩\par
\[6 5 \geq x .\]
\[Sketch the graph with equation y = f(x).\]
\textbf{(ii)} The function g is given by:\par
⎧\par
\[⎨ -x x \geq 0 ,\]
\[g(x) = 3x 0 \geq x \geq 5 ,\]
⎩\par
\[x + 10 5 \geq x .\]
Use modulus signs to write g(x) as a single formula.\par
\[\sum \sum \sum \sum\]
\textbf{(iii)} Sketch the graph with equation y = h(x), where h(x) = x 2 - x - 4 \sum x \sum + \sum x(x - 5) \sum .\par
\textbf{(iv)} The function k is given by:\par
⎧\par
\[⎨ 10x x \geq 0 ,\]
\[k(x) = 2x 2 0 \geq x \geq 5 ,\]
⎩\par
\[50 5 \geq x .\]
Use modulus signs to write k(x) as a single formula, explicitly verifying that your\par
formula is correct.\par
\end{flushleft}
\subsection*{Page 4}
\begin{flushleft}
\bigskip\textbf{Question 5.} (i) Given that a > b > c > 0 are constants, and that x, y, z are non-negative variables,\par
show that\par
\[ax + by + cz \geq a(x + y + z).\]
In the acute-angled triangle ABC, a, b and c are the lengths of sides BC, CA and AB,\par
respectively, with a > b > c. P is a point inside, or on the sides of, the triangle, and x, y\par
and z are the perpendicular distances from P to BC, CA and AB, respectively. The area of\par
the triangle is Δ.\par
\textbf{(ii)} (a) Find Δ in terms of a, b, c, x, y and z.\par
(b) Find both the minimum value of the sum of the perpendicular distances from\par
P to the three sides of the triangle and the values of x, y and z which give this\par
minimum sum, expressing your answers in terms of some or all of a, b, c and Δ.\par
\textbf{(iii)} (a) Show that, for all real a, b, c, x, y and z,\par
(a 2 + b 2 + c 2 )(x 2 + y 2 + z 2 ) = (bx - ay) 2 + (cy - bz) 2 + (az - cx) 2 + (ax + by + cz) 2 .\par
(b) Find both the minimum value of the sum of the squares of the perpendicular\par
distances from P to the three sides of the triangle and the values of x, y and z\par
which give this minimum sum, expressing your answers in terms of some or all\par
of a, b, c and Δ.\par
\textbf{(iv)} Find both the maximum value of the sum of the squares of the perpendicular distances\par
from P to the three sides of the triangle and the values of x, y and z which give this\par
maximum sum, expressing your answers in terms of some or all of a, b, c and Δ.\par
\end{flushleft}
\subsection*{Page 5}
\begin{flushleft}
\bigskip\textbf{Question 6.} In this question, you should consider only points lying in the first quadrant, that is with\par
\[x > 0 and y > 0.\]
\textbf{(i)} The equation x 2 + y 2 = 2ax defines a family of curves in the first quadrant, one curve\par
for each positive value of a. A second family of curves in the first quadrant is defined\par
2 2\par
\[by the equation x + y = 2by, where b > 0.\]
\[(a) Differentiate the equation x 2 + y 2 = 2ax implicitly with respect to x, and hence\]
show that every curve in the first family satisfies the differential equation\par
dy\par
\[2xy = y 2 - x 2 .\]
dx\par
Find similarly a differential equation, independent of b, for the second family of\par
curves.\par
(b) Hence, or otherwise, show that, at every point with y \ne x where a curve in the\par
first family meets a curve in the second family, the tangents to the two curves\par
are perpendicular.\par
A curve in the first family meets a curve in the second family at (c , c), where\par
\[c > 0. Find the equations of the tangents to the two curves at this point. Is it\]
true that where a curve in the first family meets a curve in the second family on\par
\[the line y = x, the tangents to the two curves are perpendicular?\]
\textbf{(ii)} Given the family of curves in the first quadrant y = c ln x, where c takes any non-zero\par
value, find, by solving an appropriate differential equation, a second family of curves\par
with the property that at every point where a curve in the first family meets a curve\par
in the second family, the tangents to the two curves are perpendicular.\par
\textbf{(iii)} A family of curves in the first quadrant is defined by the equation y 2 = 4k(x + k),\par
where k takes any non-zero value.\par
Show that, at every point where one curve in this family meets a second curve in the\par
family, the tangents to the two curves are perpendicular.\par
\end{flushleft}
\subsection*{Page 6}
\begin{flushleft}
\bigskip\textbf{Question 7.} Let h(z) = nz 6 + z 5 + z + n, where z is a complex number and n \leq 2 is an integer.\par
\textbf{(i)} Let w be a root of the equation h(z) = 0.\par
]\par
f(w)\par
\[(a) Show that |w 5| = , where\]
g(w)\par
\[f(z) = n\]
+ 2nRe(z) +\par
|z|2\par
\[and g(z) = n\]
2|z|2\par
+ 2nRe(z) + 1.\par
\[(b) By considering f(w) - g(w), prove by contradiction that |w| \leq 1.\]
\[(c) Show that |w| = 1.\]
\textbf{(ii)} It is given that the equation h(z) = 0 has six distinct roots, none of which is purely\par
real.\par
(a) Show that h(z) can be written in the form\par
\[h(z) = n(z 2 - a1z + 1)(z 2 - a2z + 1)(z 2 - a3z + 1),\]
\[where a^1, a^2 and a^3 are real constants.\]
\[(b) Find a^1 + a^2 + a^3 in terms of n.\]
\[(c) By considering the coefficient of z 3 in h(z), find a1a2a3 in terms of n.\]
\[(d) How many of the six roots of the equation h(z) = 0 have a negative real part?\]
Justify your answer.\par
\end{flushleft}
\subsection*{Page 7}
\begin{flushleft}
\[\geq )\]
a b\par
\bigskip\textbf{Question 8.} Let M = be a real matrix with a \ne d. The transformation represented by M has\par
c d\par
exactly two distinct invariant lines through the origin.\par
\textbf{(i)} Show that, if neither invariant line is the y-axis, then the gradients of the invariant\par
lines are the roots of the equation\par
\[bm 2 + (a - d)m - c = 0.\]
If one invariant line is the y-axis, what is the gradient of the other?\par
\textbf{(ii)} Show that, if the angle between the two invariant lines is 45 \textasciicircum{}\circ , then\par
\[(a - d) 2 = (b - c) 2 - 4bc.\]
\textbf{(iii)} Find a necessary and sufficient condition, on some or all of a, b, c and d, for the two\par
\[invariant lines to make equal angles with the line y = x.\]
\textbf{(iv)} Give an example of a matrix which satisfies both the conditions in parts (ii) and (iii).\par
\end{flushleft}
\subsection*{Page 8}
\begin{flushleft}
\bigskip\textbf{Section B: Mechanics}\par
\bigskip\textbf{Question 9.} A rectangular prism is fixed on a horizontal surface. A vertical wall, parallel to a vertical\par
face of the prism, stands at a distance d from it. A light plank, making an acute angle θ with\par
the horizontal, rests on an upper edge of the prism and is in contact with the wall below the\par
level of that edge of the prism and above the level of the horizontal plane. You may assume\par
that the plank is long enough and the prism high enough to make this possible.\par
The contact between the plank and the prism is smooth, and the coefficient of friction at\par
the contact between the plank and the wall is μ. When a heavy point mass is fixed to the\par
plank at a distance x, along the plank, from its point of contact with the wall, the system is\par
in equilibrium.\par
\textbf{(i)} Show that, if x = d sec 3 θ, then there is no frictional force acting between the plank\par
and the wall.\par
\textbf{(ii)} Show that, if x > d sec 3 θ, it is necessary that\par
x - d sec 3 θ\par
\[μ \leq\]
x tan θ\par
\[and give the corresponding inequality if x < d sec θ.\]
\textbf{(iii)} Show that\par
x sec θ\par
\[\leq .\]
d 1 + μ tan θ\par
\[Show also that, if μ < cot θ, then\]
x sec θ\par
\[\geq .\]
d 1 - μ tan θ\par
\textbf{(iv)} Show that if x is such that the point mass is fixed to the plank somewhere between\par
\[the edge of the prism and the wall, then tan θ < μ.\]
\end{flushleft}
\subsection*{Page 9}
\begin{flushleft}
\bigskip\textbf{Question 10.} (i) Show that, if a particle is projected at an angle α above the horizontal with speed u,\par
it will reach height h at a horizontal distance s from the point of projection where\par
gs\par
\[h = s tan α - .\]
\[2u2 cos2 α\]
The remainder of this question uses axes with the x- and y-axes horizontal and the z-axis\par
vertically upwards. The ground is a sloping plane with equation z = y tan θ and a road runs\par
along the x-axis. A cannon, which may have any angle of inclination and be pointed in any\par
direction, fires projectiles from ground level with speed u. Initially, the cannon is placed at\par
the origin.\par
\textbf{(ii)} Let a point P on the plane have coordinates (x , y , y tan θ). Show that the condition\par
for it to be possible for a projectile from the cannon to land at point P is\par
\[\geq )\]
2 2 4 2\par
u tan θ u sec θ\par
\[x 2 + y + \geq .\]
\[g g^2\]
\textbf{(iii)} Show that the furthest point directly up the plane that can be reached by a projectile\par
from the cannon is a distance\par
u\par
g(1 + sin θ)\par
from the cannon.\par
How far from the cannon is the furthest point directly down the plane that can be\par
reached by a projectile from it?\par
\textbf{(iv)} Find the length of road which can be reached by projectiles from the cannon.\par
The cannon is now moved to a point on the plane vertically above the y-axis, and\par
a distance r from the road. Find the value of r which maximises the length of road\par
which can be reached by projectiles from the cannon. What is this maximum length?\par
\end{flushleft}
\subsection*{Page 10}
\begin{flushleft}
\bigskip\textbf{Section C: Probability and Statistics}\par
\bigskip\textbf{Question 11.} A batch of N USB sticks is to be used on a network. Each stick has the same unknown\par
probability p of being infected with a virus. Each stick is infected, or not, independently of\par
the others.\par
The network manager decides on an integer value of T with 0 \geq T < N . If T = 0 no testing\par
takes place and the N sticks are used on the network, but if T > 0, the batch is subject to\par
the following procedure.\par
• Each of T sticks, chosen at random from the batch, undergoes a test during which it is\par
destroyed.\par
• If any of these T sticks is infected, all the remaining N - T sticks are destroyed.\par
• If none of the T sticks is infected, the remaining N - T sticks are used on the network.\par
If any stick used on the network is infected, the network has to be disinfected at a cost of\par
\pounds D, where D > 0. If no stick used on the network is infected, there is a gain of \pounds 1 for each\par
of the N - T sticks. There is no cost to testing or destroying a stick.\par
\textbf{(i)} Find an expression in terms of N , T , D and q, where q = 1 - p, for the expected net\par
loss.\par
DT\par
\textbf{(ii)} Let α = . Show that 0 \geq α < 1.\par
N (N - T + D)\par
Show that, for fixed values of N , D and T , the greatest value of the expected net loss\par
N-T\par
\[occurs when q satisfies the equation q = α.\]
D(N - T ) α k T\par
\[Show further that this greatest value is \pounds  , where k = .\]
N N - T\par
\textbf{(iii)} For fixed values of N and D, show that there is some β > 0 so that for all p < β,\par
the expression for the expected loss found in part (i) is an increasing function of T .\par
Deduce that, for small enough values of p, testing no sticks minimises the expected net\par
loss.\par
\end{flushleft}
\subsection*{Page 11}
\begin{flushleft}
\bigskip\textbf{Question 12.} The random variable X has probability density function\par
\{\par
\[kx n (1 - x) 0 \geq x \geq 1 ,\]
\[f(x) =\]
0 otherwise ,\par
where n is an integer greater than 1.\par
\textbf{(i)} Show that k = (n + 1)(n + 2) and find μ, where μ = E(X).\par
\textbf{(ii)} Show that μ is less than the median of X if\par
\[\geq )\]
n+1\par
8 2\par
\[6 - < 1 + .\]
n + 3 n + 1\par
By considering the first four terms of the expansion of the right-hand side of this\par
inequality, or otherwise, show that the median of X is greater than μ.\par
\[\geq )\]
x+1\par
\textbf{(iii)} You are given that, for positive x, 1 + is a decreasing function of x.\par
x\par
Show that the mode of X is greater than its median.\par
\end{flushleft}
\end{document}
